%% kapa: exam=1 round=33 year=2022 period=1 form=A subject=civil_law subject_ko=민법 seq=28 number=28
%% source: Q-Net 감정평가사 기출문제 (cst00309) artlSeq=5209997 file=raw/1008/5209997/2022년도 제33회 감정평가사 1차시험 1교시(민법,경제학,부동산학).pdf
%% conversion: kordoc 4.24.51 -> tools/gen1.py (auto); review=proofread-sonnet; flags=none
\begin{question}{28}
甲, 乙, 丙은 X토지를 각각 7분의 1, 7분의 2, 7분의 4의 지분으로 공유하고 있다. 이에 관한 설명으로 옳지 않은 것은? (다툼이 있으면 판례에 따름)
\begin{choices}
  \item 甲이 乙, 丙과의 협의 없이 X토지 전부를 독점적으로 점유하는 경우, 乙은 甲에 대하여 공유물의 보존행위로서 X토지의 인도를 청구할 수 없다.
  \item 丁이 X토지 전부를 불법으로 점유하는 경우, 甲은 단독으로 X토지 전부의 인도를 청구할 수 있다.
  \item 丙이 甲, 乙과의 협의 없이 X토지 전부를 戊에게 임대한 경우, 甲은 戊에게 차임 상당액의 7분의 1을 부당이득으로 반환할 것을 청구할 수 있다.
  \item 甲, 乙, 丙 사이의 X토지 사용·수익에 관한 특약이 공유지분권의 본질적 부분을 침해하지 않는 경우라면 그 특약은 丙의 특정승계인에게 승계될 수 있다.
  \item 甲은 특별한 사정이 없는 한 乙, 丙의 동의 없이 X토지에 관한 자신의 지분을 처분할 수 있다.
\end{choices}
\end{question}
\answer{3}
